\section{Thermodynamic Cost Model}

The ThermodynamicCost module injects an energy-aware conscience into the
system: every computational structure must pass an efficiency check before
it is allowed to consume resources.

\subsection{Power Measurement}

Power is sampled every second.  CPU power is estimated via psutil:
\texttt{cpu\_percent * TDP\_cpu / 100.0} (default TDP 65 W).  GPU power
is read through pynvml when an NVIDIA GPU is present.  FPGA/quantum
consumption is currently based on datasheet values but will be replaced by
direct telemetry.  The measurement is encapsulated in a dataclass:

\begin{lstlisting}[language=Python, caption={Energy snapshot}]
@dataclass(slots=True)
class EnergySnapshot:
    timestamp: float
    power_watts: float
    energy_joules: float
    temperature_c: Optional[float] = None
    cpu_percent: Optional[float] = None
    gpu_percent: Optional[float] = None
    source: EnergySource = EnergySource.UNCERTAIN
    co2_g_per_kwh: float = 0.0
\end{lstlisting}

\subsection{Efficiency Criterion}

The information value $I \in [0,1]$ of a computational structure is 
supplied by the \texttt{ComplexityProvider} interface.  In the context
of the \Apeiron{} Framework, the provider is Layer~10 (Global Coherence
Evaluation), which reports the normalised entropy of the global
coherence tensor as a proxy for informational utility.
A structure that reduces the system's internal entropy
(i.e., resolves inconsistencies or creates novel, stable regularities)
receives a high~$I$, while one that merely replicates existing knowledge
receives a low~$I$.

Given a structure $s$ that consumed $\Delta E$ Joules, it is approved iff
\[
\frac{I}{\Delta E} \ge \tau_{\mathrm{eff}},
\]
where $\tau_{\mathrm{eff}}$ is the current efficiency threshold supplied
by the AdaptiveThresholds module.  When the system is in \texttt{GREEN}
priority mode, $\tau_{\mathrm{eff}}$ is doubled.

This criterion implies that, under carbon‑constrained conditions, the
system may reject genuinely novel but computationally expensive
hypotheses in favour of cheaper, less informative ones.  The resulting
tension between environmental responsibility and cognitive completeness
is not a bug but a deliberate architectural choice: the \Core{} forces
the agent to justify every Joule it consumes, making the trade‑off
explicit and auditable.

\begin{lstlisting}[language=Python, caption={Efficiency check}]
async def evaluate(self, structure_id, computation_time, information_value):
    energy = self.monitor.measure().power_watts * computation_time
    threshold = self.adaptive.get("efficiency_threshold")
    if self.config.priority == CostPriority.GREEN:
        threshold *= 2.0
    efficiency = information_value / (energy + 1e-10)
    approved = efficiency >= threshold
    # ... log and store structure cost
    return approved, efficiency
\end{lstlisting}

\subsection{Energy Budget and Predictive Throttling}

A per-period energy budget (default 3600 J/h) is enforced.  When the budget
is exhausted, a \texttt{conserve} event is emitted, causing the Adaptive-
Thresholds module to raise all safety margins and the ChaosDetector to
increase its sensitivity.  Predictive budgeting uses an exponential moving
average of consumption per hour to forecast future usage and pre‑emptively
lower thresholds.

\subsection{Carbon Intensity and \texorpdfstring{CO$_2$}{CO2} Tracking}

Optionally, the Electricity Maps API is queried every 5 min for real‑time
carbon intensity.  If the API is unavailable, a time‑of‑day heuristic is used
(based on Netherlands grid data).  All cost decisions are logged with the
energy source and CO\(_2\) estimate, enabling transparent auditing.